% Options for packages loaded elsewhere
\PassOptionsToPackage{unicode}{hyperref}
\PassOptionsToPackage{hyphens}{url}
\documentclass[
]{article}
\usepackage{xcolor}
\usepackage{amsmath,amssymb}
\setcounter{secnumdepth}{5}
\usepackage{iftex}
\ifPDFTeX
  \usepackage[T1]{fontenc}
  \usepackage[utf8]{inputenc}
  \usepackage{textcomp} % provide euro and other symbols
\else % if luatex or xetex
\fi
\usepackage{lmodern}
\ifPDFTeX\else
  % xetex/luatex font selection
\fi
% Use upquote if available, for straight quotes in verbatim environments
\IfFileExists{microtype.sty}{% use microtype if available
  \usepackage[]{microtype}
  \UseMicrotypeSet[protrusion]{basicmath} % disable protrusion for tt fonts
}{}
\makeatletter
\@ifundefined{KOMAClassName}{% if non-KOMA class
  \IfFileExists{parskip.sty}{%
    \usepackage{parskip}
  }{% else
    \setlength{\parindent}{0pt}
    \setlength{\parskip}{6pt plus 2pt minus 1pt}}
}{% if KOMA class
  \KOMAoptions{parskip=half}}
\makeatother
\usepackage{listings}
\newcommand{\passthrough}[1]{#1}
\lstset{defaultdialect=[5.3]Lua}
\lstset{defaultdialect=[x86masm]Assembler}
\usepackage{longtable,booktabs,array}
\newcounter{none} % for unnumbered tables
\usepackage{calc} % for calculating minipage widths
% Correct order of tables after \paragraph or \subparagraph
\usepackage{etoolbox}
\makeatletter
\patchcmd\longtable{\par}{\if@noskipsec\mbox{}\fi\par}{}{}
\makeatother
\ifLuaTeX
\else
\fi
\ifLuaTeX
\fi
\setlength{\emergencystretch}{3em} % prevent overfull lines
\providecommand{\tightlist}{%
  \setlength{\itemsep}{0pt}\setlength{\parskip}{0pt}}
% ═══ 由 环境/pdf/latex/md到LaTeX.py 生成 ═══
% 中文字体：思源宋体（子集，SIL OFL 1.1）；拉丁/数学：Latin Modern（随 TeX Live 分发）
% 编译：xelatex（本仓库自带 wasm 版 XeTeX 引擎，见 环境/pdf/latex/编译LaTeX.mjs）
\usepackage[T1]{fontenc}      % 否则 ASCII 的 < > | 在 OT1 下排成 ¡ ¿ —
\usepackage{lmodern}          % 必需：数学的物理字模（否则 xdvipdfmx 拒绝出 PDF）
\usepackage{amsmath,amssymb,mathtools}
\usepackage{booktabs,longtable,array,multirow}
\usepackage{graphicx}
\usepackage{xcolor}
\usepackage{listings}
% 不用 hyperref：新版 hyperref 强制依赖 bookmark.sty，而 bundle 里没有。
% 链接命令用降级定义，保证正文里的 \href / \url 不会报错。
\providecommand{\href}[2]{#2}
\providecommand{\url}[1]{\texttt{#1}}
\providecommand{\hypertarget}[2]{#2}
\providecommand{\hyperlink}[2]{#2}
\providecommand{\urlstyle}[1]{}
\providecommand{\texorpdfstring}[2]{#1}   % hyperref 的命令，剥离后仍被模板/标题用到
\providecommand{\phantomsection}{}
\providecommand{\pdfstringdef}[2]{}
\providecommand{\hypersetup}[1]{}

% ── 三套字体：中文 / 符号 / emoji（按字符 cmap 自动选择）──
\font\zhfont="[NotoSerifSC-Regular.ttf]:script=hani" at 10.5pt
\font\symfont="[DejaVuSans.ttf]" at 10.5pt
\font\emofont="[NotoEmoji-Regular.ttf]" at 10.5pt
\font\zhmonofont="[NotoSansSC-Regular.ttf]" at 9pt
% 用法：\zh{中文} / \sym{→} / \emo{🦙}
% 注意：必须**带参数**。写成 \def\zh{{\zhfont}} 是错的——字体赋值只在小群里有效，
% 那个小群随 \zh 立刻结束，于是中文仍然用拉丁字体排，报 Missing character。
\def\zh#1{{\zhfont #1}}
\def\sym#1{{\symfont #1}}
\def\emo#1{{\emofont #1}}
% 含汉字的行内代码：listings 的 \lstinline 遇到汉字必炸（`\lst@arg ->判`
% 未定义控制字，catcode 设 12 也无效），所以直接自己排。
\def\codezh#1{{\zhmonofont #1}}

\def\zhglue{\hskip0pt plus .06em minus .01em}
\linespread{1.12}
\setlength{\parindent}{0pt}
\setlength{\parskip}{0.45em}

% ── 代码块：等宽 + 中文字体（listings 默认字体不含汉字）──
\lstset{
  basicstyle=\zhmonofont,
  breaklines=true,
  breakatwhitespace=false,
  columns=fullflexible,
  keepspaces=true,
  showstringspaces=false,
  frame=single,
  framesep=4pt,
  rulecolor=\color{gray!45},
  backgroundcolor=\color{gray!7},
  xleftmargin=6pt,
  extendedchars=true,
  tabsize=2,
  % 代码块里的 emoji：NotoSansSC 没有这些字形（listings 不做字符级回退）。
  % 用 escapeinside 开口子，转换器把 emoji 逐字包成 (*@{\emo{⭐}}@*)。
  % 试过 literate={⭐}{{\emofont ⭐}}1，会把 ASCII 字母也吞掉并报
  % `Improper alphabetic constant`，不可用。
  escapeinside={(*@}{@*)},
}
\renewcommand{\lstlistingname}{\zh{代\zhglue 码}}

% ── 表格线宽（booktabs 风格）──
\renewcommand{\arraystretch}{1.25}

% ── 标题样式 ──
\usepackage{titlesec}
\titleformat{\section}{\Large\bfseries}{\thesection}{0.6em}{}
\titlespacing*{\section}{0pt}{1.1em}{0.5em}
\urlstyle{same}
\hypersetup{
  pdftitle={\zh{课\zhglue 题}11-\zh{结\zhglue 项}},
  pdflang={zh-CN},
  hidelinks,
  pdfcreator={LaTeX via pandoc}}

\title{\zh{课\zhglue 题}11-\zh{结\zhglue 项}}
\author{}
\date{}
\begin{document}
\maketitle

{
\setcounter{tocdepth}{2}
\tableofcontents
}
\section{\zh{课\zhglue 题}11 \zh{开\zhglue 题\zhglue 简\zhglue 报} \sym{·} \zh{结\zhglue 项\zhglue （}NeurIPS \zh{报\zhglue 告} + \zh{现\zhglue 场} demo\zh{）}}\label{ux8bfeux989811-ux5f00ux9898ux7b80ux62a5-ux7ed3ux9879neurips-ux62a5ux544a-ux73b0ux573a-demo}

\begin{quote}
\zh{模\zhglue 块\zhglue ：}\textbf{M9} \zh{｜} \zh{周\zhglue 次\zhglue ：}\textbf{W15--W16\zh{（}12-21\sym{→}2027-01-03\zh{）}} \zh{｜} \zh{算\zhglue 力\zhglue ：}\textbf{\zh{弹\zhglue 性} \sym{≤}4 GPU\sym{·}\zh{小\zhglue 时}} \zh{唐\zhglue 杰\zhglue 作\zhglue 业\zhglue 对\zhglue 应\zhglue ：}\textbf{\sym{⑥} \zh{英\zhglue 文} NeurIPS \zh{格\zhglue 式} + W16 \zh{现\zhglue 场} demo}\zh{、}\textbf{\sym{⑦} \zh{作\zhglue 业} 40\% / \zh{大\zhglue 项\zhglue 目} 60\%} \zh{｜} \zh{判\zhglue 分\zhglue 来\zhglue 源\zhglue ：}\textbf{\zh{报\zhglue 告\zhglue 协\zhglue 议} + \zh{导\zhglue 师\zhglue 审\zhglue 稿} + \zh{自\zhglue 我\zhglue 攻\zhglue 击}}
\end{quote}

\subsection{\zh{一\zhglue 、\zhglue 研\zhglue 究\zhglue 背\zhglue 景}}\label{ux4e00ux7814ux7a76ux80ccux666f}

\zh{唐\zhglue 杰\zhglue 的\zhglue 原\zhglue 话\zhglue 是}''\textbf{\zh{问\zhglue 题\zhglue 、\zhglue 动\zhglue 机\zhglue 、\zhglue 方\zhglue 法\zhglue 、\zhglue 结\zhglue 果\zhglue ，\zhglue 一\zhglue 项\zhglue 都\zhglue 不\zhglue 能\zhglue 糊}}''\zh{。} \zh{这\zhglue 恰\zhglue 好\zhglue 也\zhglue 是\zhglue 本\zhglue 学\zhglue 科\zhglue 九\zhglue 份\zhglue 章\zhglue 程\zhglue 反\zhglue 复\zhglue 强\zhglue 调\zhglue 的\zhglue ：}\textbf{\zh{没\zhglue 有\zhglue 证\zhglue 据\zhglue 的\zhglue 结\zhglue 论\zhglue 不\zhglue 采\zhglue 纳\zhglue 、}no-go \zh{不\zhglue 擦\zhglue 除\zhglue 、\zhglue 误\zhglue 差\zhglue 棒\zhglue 必\zhglue 填}}\zh{。} \zh{结\zhglue 项\zhglue 不\zhglue 是}''\zh{写\zhglue 一\zhglue 篇\zhglue 总\zhglue 结}''\zh{，\zhglue 而\zhglue 是}\textbf{\zh{把} 11 \zh{个\zhglue 课\zhglue 题\zhglue 的\zhglue 证\zhglue 据\zhglue 链\zhglue 组\zhglue 装\zhglue 成\zhglue 一\zhglue 份\zhglue 可\zhglue 被\zhglue 审\zhglue 稿\zhglue 的\zhglue 作\zhglue 品}}\zh{。}

\subsection{\zh{二\zhglue 、\zhglue 研\zhglue 究\zhglue 问\zhglue 题\zhglue （\zhglue 结\zhglue 项\zhglue 版\zhglue ）}}\label{ux4e8cux7814ux7a76ux95eeux9898ux7ed3ux9879ux7248}

\begin{quote}
\textbf{\zh{我\zhglue 能\zhglue 不\zhglue 能\zhglue 用\zhglue 一\zhglue 份\zhglue 英\zhglue 文\zhglue 技\zhglue 术\zhglue 报\zhglue 告} + \zh{一\zhglue 个\zhglue 可\zhglue 运\zhglue 行\zhglue 的} demo\zh{，\zhglue 证\zhglue 明\zhglue 我\zhglue 独\zhglue 立\zhglue 完\zhglue 成\zhglue 了\zhglue 大\zhglue 模\zhglue 型\zhglue 全\zhglue 链\zhglue 路\zhglue ，\zhglue 并\zhglue 且\zhglue 知\zhglue 道\zhglue 自\zhglue 己\zhglue 的\zhglue 结\zhglue 论\zhglue 在\zhglue 哪\zhglue 里\zhglue 不\zhglue 成\zhglue 立\zhglue ？}}
\end{quote}

\subsection{\zh{三\zhglue 、\zhglue 攻\zhglue 关\zhglue 线\zhglue 索}}\label{ux4e09ux653bux5173ux7ebfux7d22}

\begin{enumerate}
\def\labelenumi{\arabic{enumi}.}
\tightlist
\item
  \textbf{\zh{报\zhglue 告\zhglue 骨\zhglue 架}}\zh{（}NeurIPS \zh{风\zhglue 格\zhglue ，}8--10 \zh{页\zhglue ）\zhglue ：}Abstract / Introduction / Related Work / Method / Experiments / Results / Discussion \& Limitations / Conclusion
\item
  \textbf{\zh{只\zhglue 写\zhglue 有\zhglue 证\zhglue 据\zhglue 的\zhglue 结\zhglue 论}}\zh{：\zhglue 每\zhglue 个\zhglue 数\zhglue 字\zhglue 都\zhglue 能\zhglue 指\zhglue 回} \passthrough{\codezh{{\zh{证}}{\zh{据}}/}} \zh{中\zhglue 的\zhglue 原\zhglue 始\zhglue 日\zhglue 志}
\item
  \textbf{Limitations \zh{必\zhglue 须\zhglue 诚\zhglue 实}}\zh{：\zhglue 规\zhglue 模\zhglue 受\zhglue 限\zhglue （}0.1B vs \zh{前\zhglue 沿\zhglue ）\zhglue 、\zhglue 算\zhglue 力\zhglue 受\zhglue 限\zhglue （\zhglue 免\zhglue 费\zhglue 档\zhglue ）\zhglue 、\zhglue 评\zhglue 测\zhglue 口\zhglue 径\zhglue 的\zhglue 局\zhglue 限}
\item
  \textbf{Demo \zh{设\zhglue 计}}\zh{：\zhglue 现\zhglue 场\zhglue 演\zhglue 示\zhglue 要}\textbf{\zh{可\zhglue 在} 3 \zh{分\zhglue 钟\zhglue 内\zhglue 跑\zhglue 完}}\zh{（\zhglue 预\zhglue 生\zhglue 成} checkpoint\zh{，\zhglue 不\zhglue 现\zhglue 场\zhglue 训\zhglue 练\zhglue ）}
\item
  \textbf{\zh{自\zhglue 我\zhglue 攻\zhglue 击} 3 \zh{条}}\zh{：\zhglue 主\zhglue 动\zhglue 写\zhglue 出}''\zh{我\zhglue 的\zhglue 结\zhglue 论\zhglue 最\zhglue 可\zhglue 能\zhglue 错\zhglue 在\zhglue 哪}''
\item
  \textbf{\zh{一\zhglue 稿\zhglue 多\zhglue 用}}\zh{：\zhglue 这\zhglue 份\zhglue 报\zhglue 告\zhglue 同\zhglue 时\zhglue 作\zhglue 为\zhglue 课\zhglue 程\zhglue 大\zhglue 论\zhglue 文\zhglue （}70\%\zh{）\zhglue 与\zhglue 唐\zhglue 杰\zhglue 作\zhglue 业}\sym{⑦}\zh{的\zhglue 交\zhglue 付\zhglue 物}
\end{enumerate}

\subsection{\zh{四\zhglue 、\zhglue 里\zhglue 程\zhglue 碑}}\label{ux56dbux91ccux7a0bux7891}

\begin{itemize}
\tightlist
\item[$\square$]
  M1 \zh{证\zhglue 据\zhglue 清\zhglue 点\zhglue ：}11 \zh{个\zhglue 课\zhglue 题\zhglue 的} \passthrough{\codezh{{\zh{证}}{\zh{据}}/}} \zh{全\zhglue 部\zhglue 齐\zhglue 备\zhglue ，\zhglue 缺\zhglue 的\zhglue 补\zhglue 齐\zhglue （\zhglue 这\zhglue 是\zhglue 一\zhglue 次}''\zh{审\zhglue 计}''\zh{）}
\item[$\square$]
  M2 \zh{英\zhglue 文\zhglue 报\zhglue 告\zhglue 初\zhglue 稿\zhglue （}Abstract + Method + Experiments \zh{先\zhglue 写\zhglue ）}
\item[$\square$]
  M3 \zh{图\zhglue 表\zhglue 规\zhglue 范\zhglue 化\zhglue ：\zhglue 所\zhglue 有\zhglue 图\zhglue 有\zhglue 坐\zhglue 标\zhglue 轴\zhglue 、\zhglue 单\zhglue 位\zhglue 、\zhglue 误\zhglue 差\zhglue 棒\zhglue 、\zhglue 样\zhglue 本\zhglue 数}
\item[$\square$]
  M4 \textbf{\zh{自\zhglue 我\zhglue 攻\zhglue 击} 3 \zh{条}} + \zh{导\zhglue 师\zhglue 审\zhglue 稿} 3--5 \zh{条} \sym{→} \zh{修\zhglue 订}
\item[$\square$]
  M5 Demo \zh{准\zhglue 备\zhglue ：\zhglue 脚\zhglue 本} + \zh{预\zhglue 生\zhglue 成\zhglue 模\zhglue 型} + 3 \zh{分\zhglue 钟\zhglue 演\zhglue 练}
\item[$\square$]
  M6 \zh{总\zhglue 复\zhglue 盘\zhglue ：\zhglue 写\zhglue 进} \passthrough{\lstinline!memory/LEARNINGS.md!}\zh{（\zhglue 哪\zhglue 些\zhglue 讲\zhglue 法\zhglue 有\zhglue 效\zhglue 、\zhglue 哪\zhglue 些\zhglue 坑\zhglue 最\zhglue 贵\zhglue ）}
\item[$\square$]
  M7 \zh{课\zhglue 程\zhglue 大\zhglue 论\zhglue 文\zhglue 定\zhglue 稿\zhglue 并\zhglue 提\zhglue 交\zhglue （\zhglue 署\zhglue 名\zhglue 直\zhglue 填\zhglue ：\zhglue 孙\zhglue 承\zhglue 泽} / 2253710052 / \zh{能\zhglue 动\zhglue 强\zhglue 基}2501\zh{）}
\end{itemize}

\subsection{\zh{五\zhglue 、\zhglue 交\zhglue 付\zhglue 物}}\label{ux4e94ux4ea4ux4ed8ux7269}

{\def\LTcaptype{none} % do not increment counter
\begin{longtable}[]{@{}lll@{}}
\toprule\noalign{}
\zh{交\zhglue 付\zhglue 物} & \zh{位\zhglue 置} & \zh{验\zhglue 收} \\
\midrule\noalign{}
\endhead
\bottomrule\noalign{}
\endlastfoot
\zh{英\zhglue 文\zhglue 报\zhglue 告} & \passthrough{\codezh{{\zh{课}}{\zh{题}}11-{\zh{结}}{\zh{项}}/report-en.md}}\zh{（}+ PDF\zh{）} & \zh{格\zhglue 式\zhglue 合\zhglue 规\zhglue 、\zhglue 结\zhglue 论\zhglue 有\zhglue 据} \\
\zh{中\zhglue 文\zhglue 大\zhglue 论\zhglue 文} & \passthrough{\codezh{{\zh{课}}{\zh{程}}{\zh{轨}}/{\zh{大}}{\zh{论}}{\zh{文}}/}} & \zh{课\zhglue 程\zhglue 要\zhglue 求\zhglue 格\zhglue 式} \\
Demo & \passthrough{\codezh{{\zh{课}}{\zh{题}}11-{\zh{结}}{\zh{项}}/{\zh{脚}}{\zh{手}}{\zh{架}}/demo.sh}} & 3 \zh{分\zhglue 钟\zhglue 内\zhglue 可\zhglue 跑} \\
\zh{复\zhglue 盘} & \passthrough{\lstinline!memory/LEARNINGS.md!} + \passthrough{\lstinline!HANDOFF.md!} & \zh{可\zhglue 被\zhglue 下\zhglue 一\zhglue 任\zhglue 直\zhglue 接\zhglue 接\zhglue 续} \\
\end{longtable}
}

\subsection{\zh{六\zhglue 、\zhglue 讲\zhglue 课\zhglue 锚\zhglue 点}}\label{ux516dux8bb2ux8bfeux951aux70b9}

\begin{itemize}
\tightlist
\item
  \textbf{\zh{讲\zhglue 义}}\zh{：}\passthrough{\codezh{M9-{\zh{前}}{\zh{沿}}.md}}\zh{（}W15 \zh{展\zhglue 开\zhglue ）}
\item
  \textbf{\zh{对\zhglue 标}}\zh{：}\textbf{rasbt ch3 \zh{的}''\zh{先\zhglue 建\zhglue 评\zhglue 测}''\zh{精\zhglue 神}}\zh{贯\zhglue 穿\zhglue 全\zhglue 文} \sym{·} CS336 \zh{报\zhglue 告\zhglue 规\zhglue 范} \sym{·} \passthrough{\codezh{{\zh{论}}{\zh{文}}{\zh{精}}{\zh{读}}/}} \zh{的\zhglue 英\zhglue 文\zhglue 写\zhglue 作\zhglue 规\zhglue 范}
\item
  \textbf{\zh{搭\zhglue 桥}}\zh{：}\passthrough{\lstinline!turbine-blade-ai-platform!}\zh{（\zhglue 应\zhglue 用\zhglue 章\zhglue 节\zhglue ）\zhglue ；}\passthrough{\codezh{{\zh{人}}{\zh{工}}{\zh{智}}{\zh{能}}{\zh{基}}{\zh{础}}{\zh{与}}{\zh{应}}{\zh{用}}/{\zh{课}}{\zh{程}}{\zh{轨}}/{\zh{大}}{\zh{论}}{\zh{文}}}}\zh{（\zhglue 一\zhglue 稿\zhglue 两\zhglue 用\zhglue ）}
\end{itemize}

\subsection{\zh{七\zhglue 、\zhglue 自\zhglue 检\zhglue 清\zhglue 单\zhglue （\zhglue 结\zhglue 项\zhglue 前\zhglue 必\zhglue 过\zhglue ，\zhglue 抄\zhglue 自\zhglue 《\zhglue 判\zhglue 分\zhglue 与\zhglue 验\zhglue 收\zhglue 标\zhglue 准\zhglue 》}\sym{§}\zh{五\zhglue ）}}\label{ux4e03ux81eaux68c0ux6e05ux5355ux7ed3ux9879ux524dux5fc5ux8fc7ux6284ux81eaux5224ux5206ux4e0eux9a8cux6536ux6807ux51c6ux4e94}

\begin{enumerate}
\def\labelenumi{\arabic{enumi}.}
\tightlist
\item
  \zh{每\zhglue 个\zhglue 数\zhglue 字\zhglue 都\zhglue 能\zhglue 在} \passthrough{\codezh{{\zh{证}}{\zh{据}}/}} \zh{找\zhglue 到\zhglue 原\zhglue 始\zhglue 日\zhglue 志\zhglue 吗\zhglue ？}
\item
  ``\zh{变\zhglue 好\zhglue 了}''\zh{有\zhglue 没\zhglue 有\zhglue 可\zhglue 能\zhglue 是\zhglue 口\zhglue 径\zhglue 变\zhglue 了} / \zh{只\zhglue 挑\zhglue 了\zhglue 好\zhglue 种\zhglue 子\zhglue ？}
\item
  \zh{有\zhglue 没\zhglue 有\zhglue 把}\textbf{\zh{训\zhglue 练\zhglue 集\zhglue 表\zhglue 现}}\zh{当\zhglue 能\zhglue 力\zhglue ？}
\item
  \zh{每\zhglue 个}''\zh{因\zhglue 为}\ldots \zh{所\zhglue 以}\ldots``\zh{有\zhglue 没\zhglue 有\zhglue 反\zhglue 事\zhglue 实\zhglue 检\zhglue 验\zhglue ？}
\item
  \zh{图\zhglue 表\zhglue 的\zhglue 坐\zhglue 标\zhglue 轴\zhglue 与\zhglue 误\zhglue 差\zhglue 棒\zhglue 都\zhglue 标\zhglue 了\zhglue 吗\zhglue ？}
\item
  \zh{我\zhglue 写\zhglue 下\zhglue 的\zhglue 最\zhglue 强\zhglue 反\zhglue 方\zhglue 论\zhglue 点\zhglue 是\zhglue 什\zhglue 么\zhglue ？}
\end{enumerate}

\subsection{\zh{八\zhglue 、\zhglue 风\zhglue 险\zhglue 与\zhglue 提\zhglue 醒}}\label{ux516bux98ceux9669ux4e0eux63d0ux9192}

\begin{itemize}
\tightlist
\item
  \sym{⚠️} \zh{报\zhglue 告\zhglue 最\zhglue 易\zhglue 犯\zhglue 的\zhglue 错\zhglue 是}\textbf{\zh{夸\zhglue 大}}\zh{：\zhglue 规\zhglue 模\zhglue 受\zhglue 限\zhglue 就\zhglue 要\zhglue 说\zhglue 受\zhglue 限\zhglue ，\zhglue 这\zhglue 反\zhglue 而\zhglue 是\zhglue 可\zhglue 信\zhglue 度\zhglue 来\zhglue 源\zhglue 。}
\item
  \sym{⚠️} \zh{不\zhglue 要\zhglue 为\zhglue 了}''\zh{好\zhglue 看}''\zh{删\zhglue 掉} no-go \zh{记\zhglue 录}------\zh{用\zhglue 户\zhglue 明\zhglue 确\zhglue 要\zhglue 求}''\zh{负\zhglue 结\zhglue 果\zhglue 不\zhglue 许\zhglue 擦\zhglue 除}''\zh{。}
\item
  \sym{⚠️} W16 \zh{现\zhglue 场} demo \zh{必\zhglue 须\zhglue 有}\textbf{\zh{失\zhglue 败\zhglue 预\zhglue 案}}\zh{（\zhglue 预\zhglue 生\zhglue 成\zhglue 结\zhglue 果} + \zh{录\zhglue 屏\zhglue ）\zhglue ，\zhglue 不\zhglue 依\zhglue 赖\zhglue 现\zhglue 场\zhglue 网\zhglue 络\zhglue 与\zhglue 算\zhglue 力\zhglue 。}
\end{itemize}

\end{document}